% Generated draft: requires mathematical and source review.
\documentclass[11pt]{article}
\usepackage[margin=25mm]{geometry}
\usepackage{fontspec}
\setmainfont{Latin Modern Roman}
\usepackage{amsmath,amssymb,mathtools,unicode-math}
\setmathfont{Latin Modern Math}
\usepackage{enumitem,needspace,xurl,hyperref}
\setlength{\parindent}{0pt}
\setlength{\parskip}{5pt}
\setlength{\emergencystretch}{3em}
\newcommand{\N}{\mathbb{N}}
\newcommand{\Z}{\mathbb{Z}}
\newcommand{\Q}{\mathbb{Q}}
\newcommand{\R}{\mathbb{R}}
\newcommand{\C}{\mathbb{C}}
\newcommand{\F}{\mathbb{F}}
\newcommand{\A}{\mathbb{A}}
\newcommand{\PP}{\mathbb{P}}
\newcommand{\E}{\mathbb{E}}
\newcommand{\eps}{\varepsilon}
\newcommand{\dd}{\,\mathrm{d}}
\begin{document}
\section*{UnsolvedMath ID: 801041}
\textbf{Source problem code:} catalogue entry
\subsection{Definitions and mathematical statement}
Interpret the question in the usual setting of Borel actions on standard Borel spaces. A standard Borel space is a measurable space isomorphic to the Borel measurable space of a Polish space. This ambient-space hypothesis is conventional but is not explicit in the supplied statement.

Let \(G\) be a countable group, equipped with the discrete Borel structure. Say that \(G\) is amenable if, for every finite set \(K\subseteq G\) and every \(\varepsilon>0\), there is a nonempty finite set \(F\subseteq G\) such that
\[
 |kF\mathbin{\triangle}F|<\varepsilon |F|\qquad\text{for every }k\in K,
\]
where \(\triangle\) denotes symmetric difference. A Borel action of \(G\) on a standard Borel space \(X\) is a Borel map \(G\times X\to X\), denoted \((g,x)\mapsto g\cdot x\), satisfying \(e\cdot x=x\) and \((gh)\cdot x=g\cdot(h\cdot x)\). Its orbit equivalence relation is
\[
 E_G^X=\{(x,y)\in X\times X:\exists g\in G\ (y=g\cdot x)\}.
\]
A Borel equivalence relation \(E\) on \(X\) is hyperfinite if there are Borel equivalence relations \(E_n\subseteq X\times X\), indexed by integers \(n\geq 0\), such that every \(E_n\)-class is finite and
\[
 E_n\subseteq E_{n+1}\quad(n\geq0),\qquad E=\bigcup_{n\geq0}E_n.
\]
The finite classes need not have a common bound on their sizes.

\textbf{Question (attributed in the supplied record to Weiss).} For every countable amenable group \(G\), every standard Borel space \(X\), and every Borel action \(G\curvearrowright X\), is \(E_G^X\) hyperfinite?

No freeness assumption or probability measure is imposed. Hyperfiniteness here is required on all of \(X\), without discarding an exceptional set. These specifications express the conventional Borel interpretation; the supplied one-sentence statement does not explicitly specify the ambient measurable space or distinguish global hyperfiniteness from measure-dependent variants.
\subsection{Short English statement}
Weiss' question asks whether every Borel action of a countable amenable group on a standard Borel space has an orbit relation that is an increasing union of Borel equivalence relations with finite classes. The standard Borel setting is a conventional interpretation of the supplied wording and requires review. The registry labels the problem open; that status is supplied metadata, not an independently verified assessment.
\subsection{Sources}
\begin{itemize}
\item Local UnsolvedMath registry, permanent ID \(801041\), exact source problem code \texttt{catalogue entry}; supplied title: \emph{Open Problems in Mathematical Logic --- Problem 53 (Weiss' question)}. The supplied record attributes the question to Weiss and lists Barmpalias, Gao et al. as the source-list authors.
\item Barmpalias, Gao et al., \emph{Open Problems in Mathematical Logic} (2026), Problem 53. Supplied source URL: \url{https://arxiv.org/abs/2608.26628}. The supplied provenance records access on 2026-10-16 and extraction by \texttt{arxiv\_latex}; neither access nor the source contents have been independently verified for this draft.
\item The local record supplies the status \texttt{open} and upstream revision \texttt{d371ee917ac9647fbb06cd6fc529deec6f3dda82}. These are source metadata, not independent mathematical or literature verification.
\end{itemize}
\end{document}
